\chapter{شرط‌ها و تصمیم‌گیری}
\label{ch:conditionals}

تا اینجا برنامه‌های ما خط‌به‌خط و بدونِ هیچ انتخابی پیش می‌رفتند. اما زندگیِ
واقعی پر از تصمیم است: «اگر باران آمد، چتر بردار»، «اگر نمره از ۱۰ بیشتر شد،
قبول هستی». در این فصل یاد می‌گیریم که برنامه‌مان هم بتواند تصمیم بگیرد.

\section{اگر: نخستین تصمیم}

ساده‌ترین شکلِ تصمیم‌گیری با واژهٔ \code{اگر} (معادلِ \code{if}) انجام می‌شود:

\begin{salam}
تابع اصلی:
    سن : صحیح = 20
    اگر سن >= 18:
        چاپ "شما بزرگسال هستید"
    پایان
پایان
\end{salam}

این برنامه را چنین بخوانید: «\emph{اگر} سن بزرگ‌تر یا مساویِ ۱۸ بود، آنگاه
پیام را چاپ کن.» اگر شرط \emph{درست} باشد، بدنهٔ \code{اگر} اجرا می‌شود؛ وگرنه
سلام از رویِ آن می‌پرد.

ساختار همیشه یکسان است:

\begin{center}
\code{اگر} \;\textit{شرط}\; \code{:} \quad\ldots\;بدنه\;\ldots\quad \code{پایان}
\end{center}

شرط، هر عبارتی است که به \code{درست} یا \code{نادرست} برسد درست همان
عبارت‌های مقایسه‌ای و منطقیِ فصلِ پیش.

\section{وگرنه: راهِ دوم}

گاهی می‌خواهیم اگر شرط درست \emph{نبود}، کارِ دیگری انجام شود. برای این از
\code{وگرنه} (معادلِ \code{else}) استفاده می‌کنیم:

\begin{salam}
تابع اصلی:
    سن : صحیح = 15
    اگر سن >= 18:
        چاپ "بزرگسال"
    وگرنه:
        چاپ "نوجوان یا کودک"
    پایان
پایان
\end{salam}

اکنون دقیقاً یکی از دو پیام چاپ می‌شود: یا «بزرگسال» (اگر شرط درست باشد) یا
«نوجوان یا کودک» (اگر نادرست باشد). هرگز هر دو با هم چاپ نمی‌شوند.

\section{وگرنه اگر: چند راه}

زندگی همیشه دوراهی نیست؛ گاهی چند گزینه داریم. برای این حالت، \code{وگرنه}
را به‌همراهِ شرطِ بعدی (معادلِ \code{else if}) پشتِ سرِ هم می‌چینیم. به این
مثالِ واقعی که نمره را به طبقهٔ حرفی تبدیل می‌کند نگاه کنید:

\begin{salam}
تابع طبقه(ن : صحیح) : رشته:
    اگر ن >= 90:
        بازگشت "الف"
    وگرنه ن >= 70:
        بازگشت "ب"
    وگرنه:
        بازگشت "ج"
    پایان
پایان

تابع اصلی:
    چاپ طبقه(95)
    چاپ طبقه(75)
    چاپ طبقه(40)
پایان
\end{salam}

\begin{tip}
دقت کنید که سلام واژهٔ \code{اگر} را دوباره بعد از \code{وگرنه} تکرار
نمی‌کند. اگر بی‌درنگ بعد از \code{وگرنه} یک شرط بیاید، همان نقشِ
\code{else if} را دارد؛ و اگر \code{وگرنه} بدونِ هیچ شرطی بیاید، همان
شاخهٔ نهایی و پیش‌فرض است.
\end{tip}

خروجی:

\begin{salamen}[language={}]
الف
ب
ج
\end{salamen}

سلام شرط‌ها را \emph{از بالا به پایین} می‌آزماید و به \emph{نخستین} شرطِ درست
که رسید، بدنهٔ آن را اجرا می‌کند و باقی را نادیده می‌گیرد. پس برای نمرهٔ ۹۵،
نخستین شرط (\code{>= 90}) درست است و «الف» بازگردانده می‌شود؛ سلام دیگر شرط‌های
بعدی را نمی‌آزماید.

\begin{warn}
ترتیبِ شرط‌ها اهمیت دارد! اگر در مثالِ بالا شرطِ \code{ن >= 70} را پیش از
\code{ن >= 90} می‌نوشتیم، نمرهٔ ۹۵ نیز در شرطِ ۷۰ گیر می‌افتاد و به‌اشتباه «ب»
می‌گرفت. همیشه شرط‌های \emph{سخت‌گیرانه‌تر} (محدودتر) را زودتر بیاورید.
\end{warn}

\begin{note}
در این مثال از \code{بازگشت} درونِ یک تابع استفاده کردیم تا نتیجه را برگردانیم.
\code{بازگشت} هم مقدار را به بیرونِ تابع می‌فرستد و هم بی‌درنگ از تابع خارج
می‌شود. توابع را در فصلِ ۶ به‌طورِ کامل می‌شناسیم.
\end{note}

\section{شرط در یک خط}

وقتی بدنهٔ شرط کوتاه است، می‌توانید همه را در یک خط بنویسید. این مثالِ واقعیِ
«سالِ کبیسه» سه شیوهٔ فشرده را نشان می‌دهد:

\begin{salam}
// سال کبیسه
تابع کبیسه_است(سال : صحیح) : منطقی:
    اگر سال % 400 == 0: بازگشت درست پایان
    اگر سال % 100 == 0: بازگشت نادرست پایان
    بازگشت سال % 4 == 0
پایان

تابع اصلی:
    متغیر سالها : صحیح[5] = [1900, 2000, 2012, 2021, 2024]
    برای متغیر i = 0, i < 5, i = i + 1:
        اگر کبیسه_است(سالها[i]): چاپ سالها[i], "کبیسه است"
        وگرنه:                    چاپ سالها[i], "کبیسه نیست"
        پایان
    پایان
پایان
\end{salam}

اینجا چند چیزِ تازه می‌بینیم:
\begin{itemize}
  \item شرط و بدنه و \code{پایان} در یک خط: \code{اگر سال \% 400 == 0: بازگشت
        درست پایان}. این شکلِ فشرده برای شرط‌های کوتاه بسیار خواناست.
  \item تابعی که یک مقدارِ \code{منطقی} برمی‌گرداند خودِ
        \code{سال \% 4 == 0} یک عبارتِ منطقی است و مستقیم بازگردانده می‌شود.
  \item یک آرایه و یک حلقه که در فصل‌های بعد به آن‌ها می‌رسیم.
\end{itemize}

\begin{deep}[نگاهی به درون: منطقِ سالِ کبیسه]
چرا قاعدهٔ کبیسه این‌قدر پیچیده است؟ سالی کبیسه است که بر ۴ بخش‌پذیر باشد ---
\emph{مگر} آنکه بر ۱۰۰ بخش‌پذیر باشد (که آنگاه کبیسه نیست)، \emph{مگر} آنکه بر
۴۰۰ بخش‌پذیر باشد (که باز کبیسه است). به همین دلیل کد، نخست ۴۰۰، سپس ۱۰۰ و در
آخر ۴ را می‌آزماید ترتیبی که دقیقاً منطقِ این استثناها را پیاده می‌کند.
\end{deep}

\section{شرط‌های تودرتو}

می‌توانید یک \code{اگر} را \emph{درونِ} یک \code{اگر} دیگر بگذارید. به این
نمونه دقت کنید که نخست بررسی می‌کند آیا کاربر وارد شده، و سپس آیا مدیر است:

\begin{salam}
تابع اصلی:
    واردشده : منطقی = درست
    مدیر : منطقی = نادرست
    اگر واردشده:
        اگر مدیر:
            چاپ "خوش آمدید، مدیرِ گرامی"
        وگرنه:
            چاپ "خوش آمدید، کاربرِ عزیز"
        پایان
    وگرنه:
        چاپ "لطفاً ابتدا وارد شوید"
    پایان
پایان
\end{salam}

\begin{tip}
شرط‌های تودرتوی عمیق، کد را سخت‌خوان می‌کنند. اغلب می‌توان به‌جای تودرتو کردن،
شرط‌ها را با \code{\&\&} ترکیب کرد. مثلاً
\code{اگر واردشده \&\& مدیر:} کوتاه‌تر و خواناتر از دو \code{اگر}ِ تودرتوست ---
البته اگر به آن «وگرنه»ی میانی نیاز نداشته باشید.
\end{tip}

\section{مثالِ کامل: بازیِ حدسِ ساده}

بیایید همهٔ آموخته‌ها را در یک مثالِ کوچک گرد آوریم. این برنامه یک دما را
ارزیابی می‌کند:

\begin{salam}
تابع وضعیت_هوا(دما : صحیح) : رشته:
    اگر دما < 0:
        بازگشت "یخبندان"
    وگرنه دما < 15:
        بازگشت "سرد"
    وگرنه دما < 30:
        بازگشت "معتدل"
    وگرنه:
        بازگشت "گرم"
    پایان
پایان

تابع اصلی:
    چاپ "۵- درجه:", وضعیت_هوا(-5)
    چاپ "۱۰ درجه:", وضعیت_هوا(10)
    چاپ "۲۵ درجه:", وضعیت_هوا(25)
    چاپ "۳۸ درجه:", وضعیت_هوا(38)
پایان
\end{salam}

این الگو یک تابع که با زنجیره‌ای از \code{وگرنه اگر} یک مقدار را به یک
دسته نگاشت می‌کند یکی از پرکاربردترین الگوهای برنامه‌نویسی است.

\section{خلاصهٔ فصل}

\begin{itemize}
  \item \code{اگر} بدنه را تنها وقتی اجرا می‌کند که شرط درست باشد.
  \item \code{وگرنه} راهِ جایگزین را برای حالتِ نادرست فراهم می‌کند.
  \item \code{وگرنه اگر} چند گزینه را پشتِ سرِ هم می‌آزماید.
  \item شرط‌ها از بالا به پایین آزموده می‌شوند و نخستین درست، برنده است.
  \item شرط‌های کوتاه را می‌توان در یک خط نوشت.
  \item به‌جای تودرتوی عمیق، شرط‌ها را با \code{\&\&} و \code{||} ترکیب کنید.
\end{itemize}

\section{تمرین‌ها}

\exercise برنامه‌ای بنویسید که یک عدد بگیرد و بگوید مثبت، منفی یا صفر است.

\exercise تابعی به نامِ \code{بزرگ‌تر} بنویسید که دو عدد بگیرد و بزرگ‌ترشان را
برگرداند (با \code{اگر} و \code{وگرنه}).

\exercise برنامه‌ای بنویسید که نمره‌ای بین ۰ تا ۲۰ بگیرد و بگوید: عالی (۱۸ به
بالا)، خوب (۱۵ تا ۱۸)، قبول (۱۰ تا ۱۵) یا مردود (زیرِ ۱۰).

\exercise با کمکِ عملگرِ \code{\%}، تابعی بنویسید که بگوید عددی زوج است یا فرد.

\exercise تابعِ \code{کبیسه\_است} را بازنویسی کنید، اما این بار به‌جای سه
\code{اگر}ِ جداگانه، همه را در \emph{یک} عبارتِ منطقیِ بزرگ بنویسید (راهنمایی:
از \code{\&\&}، \code{||} و پرانتز استفاده کنید).
